\documentclass[../../main.tex]{subfiles}
\begin{document}

{\bf [解答]}

漸化式から $k \in \mathbb{Z}_{\ge 0}$ として
\begin{align}
  f^{4k}(x)   = \sin x, \quad
  f^{4k+1}(x) = \cos x, \quad
  f^{4k+2}(x) = -\sin x, \quad
  f^{4k+3}(x) = -\cos x \label{eq:1}
\end{align}
である．以下点 $P$ の $x$ 座標を $p$ として，$n$を$4$で割ったあまりで場合分けする．
それぞれの場合で$t_1, t_2$ の接線ベクトルの内積が$0$となれば良い．

\subparagraph{$1^\circ$ $n=4k+1$の時}

この時は\cref{eq:1}より
\begin{align}
  C_1 : y = x\sin x, \quad
  C_2 : y = \cos x
\end{align}
だから，点$P$でこれらが交点を持つことより
\begin{align}
  p \sin p = \cos p \label{eq:2}
\end{align}
となる．
$t_1, t_2$ の接線ベクトルは $
\begin{pmatrix} 1 \\ p \cos p + \sin p
\end{pmatrix}$, $
\begin{pmatrix} 1 \\ -\sin p
\end{pmatrix}$
だからこれらの内積は
\begin{align}
  \begin{pmatrix} 1 \\ p \cos p + \sin p
  \end{pmatrix} \cdot
  \begin{pmatrix} 1 \\ -\sin p
  \end{pmatrix} = \cos p (\cos p - p \sin p) = 0
\end{align}
と$0$になる，つまり$t_1$ と $t_2$ は直交する．

\subparagraph{$2^\circ$ $n=4k+2$の時}

この時\cref{eq:1}より
\begin{align}
  C_1 : y = x\cos x, \quad
  C_2 : y = -\sin x
\end{align}
だから，$p \cos p = -\sin p$ であり，接線ベクトルは
$
\begin{pmatrix} 1 \\ \cos p - p \sin p
\end{pmatrix}$, $
\begin{pmatrix} 1 \\ -\cos p
\end{pmatrix}$
となり，$1^\circ$ と同じく $t_1, t_2$ は直交する．

\subparagraph{$3^\circ$ $n=4k+3$の時}

この時\cref{eq:1}より
\begin{align}
  C_1 : y = -x\sin x, \quad
  C_2 : y = -\cos x
\end{align}
だから$-p \sin p = -\cos p$ であり，接線ベクトルは
$
\begin{pmatrix} 1 \\ -p \cos p - \sin p
\end{pmatrix}$, $
\begin{pmatrix} 1 \\ \sin p
\end{pmatrix}$
だから，$1^\circ$ と同じく $t_1, t_2$ は直交する．

\subparagraph{$4^\circ$ $n=4k$の時}

この時\cref{eq:1}より
\begin{align}
  C_1 : y = -x\cos x,\quad
  C_2 : y = \sin x
\end{align}
だから$-p \cos p = \sin p$ であり，接線ベクトルは
$
\begin{pmatrix} 1 \\ -\cos p + p \sin p
\end{pmatrix}$, $
\begin{pmatrix} 1 \\ \cos p
\end{pmatrix}$
だから，$1^\circ$ と同じく $t_1, t_2$ は直交する．
\casesend

以上からいずれの場合も$t_1,t_2$は直交する．
故に題意は示された．

\newpage

\end{document}
